\documentclass[10pt,a4paper]{amsart}
\usepackage[margin=19mm]{geometry}
\usepackage{amsmath,amssymb,mathtools}
\usepackage[scr=rsfs,cal=euler]{mathalfa}
\usepackage{xfrac}
\usepackage[unicode,hidelinks,bookmarks=false]{hyperref}
\newcommand{\ie}{i.e.\ }
\newcommand{\cf}{cf.\ }
\newcommand{\resp}{respectively\ }
\newcommand{\Subsection}{\subsection}
\providecommand{\nobreakdash}{\nobreak\hbox{-}}
\setlength{\emergencystretch}{2em}
\multlinegap0pt
\usepackage{amssymb}
\usepackage{enumerate}
\newtheorem{lemma}{Lemma}
\newtheorem{theorem}{Theorem}
\title{Bergeron's Conjecture \& A Tale of Two Binomial Coefficients}
\author{Tewodros Amdeberhan, Matthias Beck}
\date{Source: arXiv:2607.04050v1 (CC BY 4.0)}
\begin{document}
\maketitle

\noindent\emph{Source and licence.} Bergeron's Conjecture \& A Tale of Two Binomial Coefficients. Version arXiv:2607.04050v1 (CC BY 4.0), distributed by arXiv under the Creative Commons Attribution 4.0 International licence, http://creativecommons.org/licenses/by/4.0/, as recorded in the arXiv licence metadata for this exact version and shown on the abstract page https://arxiv.org/abs/2607.04050v1. Full licence notice: \texttt{licenses/arxiv-2607.04050-CC-BY-4.0.txt}.\par\smallskip

\noindent\textit{Editorial scope.} Counted unit: the article's theorem theorem (source0.tex lines 385--388). The article's complete original proof of it is delivered with it (source0.tex lines 390--459, one \texttt{proof} environment of 70 lines). No mathematics is written, repaired or re-derived: the statement and the proof are the article's own lines, in the article's own notation, and no retained line is substituted or re-flowed.  Retained uncounted as the source-local context the proof needs: lines 316--319.  Nothing was omitted from any retained span, so the package carries no omission record.  No retained line contains a citation, so the wrapper carries no bibliography and no bibliography entry.  No retained line contains a figure environment, an image-inclusion command or drawing code, so no image is delivered and the package declares no asset dependency.  Editorial formatting only, in addition: an AMS \texttt{amsart} preamble that restates the article's own declarations and macros used by the retained lines, namely the environments \texttt{lemma}, \texttt{theorem} on separate counters, together with the standard packages those lines need.  The retained environments print in this document as Lemma 1, Theorem 1 in order of appearance, whereas the article prints the same items with the numbering of its own full document.  The complete native source of the article as distributed in the arXiv e-print for this exact version is bundled unchanged as \texttt{sources/204444/source0.tex}.
\begin{lemma} \label{lemma1}
For integers $b> a \ge 1$,
$$ \binom{b+2a}b \ge \binom{a+2b}{a}.$$
\end{lemma}

\begin{theorem}
For integers $b> a \ge 1$ and integer $\beta\geq2$, 
$$\binom{b+\beta a}b \ge \binom{a+\beta b}{a}.$$
\end{theorem}

\begin{proof} For convenience, we write $b=na+r$ where $n\geq1$ and $a>r\geq0$ so that the assertion is tantamount to
$$ \binom{(n+\beta)a+r}{\beta a} \ge \binom{(\beta n+1)a+\beta r}{a}.$$
 Let $U$ be a set with $|U| = (n+\beta)a+r$, and let $V$ be a disjoint set with $|V| = (\beta-1)(n-1)a+(\beta-1)r$. Define $S = U \cup V$, so $|S| = (\beta n+1)a+\beta r$. 

Let's represent the two binomial coefficients as collections of subsets:

\begin{itemize}
    \item Let $\mathcal{A}$ be the collection of all $a$-element subsets of $S$, i.e., $\mathcal{A} = \{ A \subseteq S : |A| = a \}$. Then
    $$
    |\mathcal{A}|=\binom{(\beta n+1)a+\beta r}{a}.
    $$

    \item Let $\mathcal{B}$ be the collection of all $\beta a$-element subsets of $U$, i.e., $\mathcal B = \{ B \subseteq U : |B| = \beta a \}$. Then
    $$
    |\mathcal{B}|=\binom{(n+\beta)a+r}{\beta a}.
    $$
\end{itemize}
We say that $A\in\mathcal{A}$ is \emph{linked to} $B\in\mathcal{B}$ if $A\cap U \subseteq B$.
Let $E$ denote the total number of such links.

We proceed by induction on $\beta\geq2$. The base case $\beta=2$ is exactly the content of Lemma~\ref{lemma1}. 
% Assume $P(\beta-1)$ holds true and we shall prove $P(\beta)$ is valid as well.
For the induction step, assume the inequality in the statement of our theorem holds for $\beta - 1$.

\subsection*{Counting from the $\mathcal{B}$-side.}

Fix a subset $B\in\mathcal{B}$. For a subset $A\in\mathcal{A}$ to satisfy $A\cap U\subseteq B$, the set $A$ cannot contain any element of $U\setminus B$. Hence all elements of $A$ must be chosen from the disjoint union $B\cup V$. Since $|B|=\beta a$ and $|V|=(\beta-1)((n-1)a+r)$, we have $|B\cup V|=((\beta-1)n+1)a+(\beta-1)r$.

Therefore the number of subsets $A$ linked to a fixed $B$ equals
$$
\binom{((\beta-1)n+1)a+(\beta-1)r}a.
$$
Because this holds for every $B\in\mathcal{B}$, we gather that
\begin{align*}
E  \ &= \ |\mathcal{B}|\cdot \binom{((\beta-1)n+1)a+(\beta-1)r}a \\ 
&= \ \binom{(n+\beta)a+r}{\beta a} \binom{((\beta-1)n+1)a+(\beta-1)r}a \, .
\end{align*}

\subsection*{Counting from the $\mathcal{A}$-side.}
Fix a subset $A\in\mathcal{A}$ and let $k=|A\cap U|$. Since $|A|=a$, it must be that $0\leq k\leq a$.
To construct a subset $B\in\mathcal{B}$ linked to $A$, the set $B$ must contain all $k$ elements of $A\cap U$. 
The remaining $\beta a-k$ elements of $B$ must then be chosen from the remaining $(n+\beta)a+r-k$ members of $U$.
Hence the number of such subsets $B$ equals
$$\binom{(n+\beta)a+r-k}{\beta a-k}=\binom{(n+\beta)a+r-k}{na+r},$$
using symmetry of the binomials. On the other hand, we have $(n+\beta)a+r-k \geq (n+\beta-1)a+r$
based on $k\leq a$. Using this observation together with the monotonicity of the binomials $\binom{n}k\geq\binom{m}k$, when $n\geq m$, it follows that
$$\binom{(n+\beta)a+r-k}{na+r} \geq \binom{(n+\beta-1)a+r}{na+r}=\binom{(n+\beta-1)a+r}{(\beta-1)a}.$$

Thus every subset $A\in\mathcal{A}$ links to at least
$$
\binom{(n+\beta-1)a+r}{(\beta-1)a}
$$
subsets in $\mathcal{B}$. Consequently,
\begin{align*}
 E &\geq |\mathcal{A}|\cdot \binom{(n+\beta-1)a+r}{(\beta-1)a} = \binom{(\beta n+1)a+\beta r}a \binom{(n+\beta-1)a+r}{(\beta-1)a} \\
& \geq \binom{(\beta n+1)a+\beta r}a \binom{((\beta-1)n+1)a+(\beta-1)r}a
\end{align*}
where the last inequality is due to the induction hypothesis. % $P(\beta-1)$. 
Combining the two counts for $E$, from the $\mathcal{B}$-side and from $\mathcal{A}$-side, 
we arrive at the estimate
\begin{align*}
&\binom{(n+\beta)a+r}{\beta a} \binom{((\beta-1)n+1)a+(\beta-1)r}a \\
&\geq \binom{(\beta n+1)a+\beta r}a \binom{((\beta-1)n+1)a+(\beta-1)r}a \, .
\end{align*}
Canceling the positive factor $\binom{((\beta-1)n+1)a+(\beta-1)r}a$ from both gives
$$
\binom{(n+\beta)a+r}{\beta a} \ge \binom{(\beta n+1)a+\beta r}a
$$
and hence the proof follows.
\end{proof}

\end{document}
